\documentclass[11pt,a4paper]{article}
\usepackage[utf8]{inputenc}
\usepackage[T1]{fontenc}
\usepackage[spanish]{babel}
\usepackage{geometry}
\usepackage{amsmath}
\usepackage{amssymb}
\usepackage{enumitem}
\usepackage{booktabs}
\usepackage{fancyhdr}
\usepackage{xcolor}
\usepackage{mdframed}
\usepackage{titlesec}
\usepackage{multicol}

\geometry{top=2cm, bottom=2.5cm, left=2.5cm, right=2.5cm}

\pagestyle{fancy}
\fancyhf{}
\rhead{Tecnolog\'ia Digital VI --- 1er Semestre 2026}
\lhead{Banco de Preguntas Te\'oricas}
\rfoot{P\'agina \thepage}

\definecolor{tdblue}{RGB}{0,74,135}
\definecolor{tdgreen}{RGB}{0,112,60}
\definecolor{tdred}{RGB}{180,35,24}
\definecolor{tdorange}{RGB}{200,100,0}
\definecolor{lightgray}{RGB}{245,245,245}
\definecolor{lightblue}{RGB}{230,240,255}
\definecolor{lightgreen}{RGB}{230,248,235}
\definecolor{lightyellow}{RGB}{255,252,220}

% Caja de sección
\newmdenv[linecolor=tdblue, linewidth=2pt, roundcorner=5pt,
  backgroundcolor=lightblue, innertopmargin=8pt, innerbottommargin=8pt,
  innerleftmargin=10pt]{sectionbox}

% Caja de advertencia (pregunta clásica)
\newmdenv[linecolor=tdred, linewidth=2pt, roundcorner=5pt,
  backgroundcolor=lightyellow, innertopmargin=8pt, innerbottommargin=8pt,
  innerleftmargin=10pt]{classicbox}

% Caja de respuesta
\newmdenv[linecolor=tdgreen, linewidth=1pt, roundcorner=3pt,
  backgroundcolor=lightgreen, innertopmargin=6pt, innerbottommargin=6pt,
  innerleftmargin=8pt]{answerbox}

\titleformat{\section}{\large\bfseries\color{tdblue}}{}{0em}{}[\titlerule]
\titleformat{\subsection}{\normalsize\bfseries\color{tdblue}}{}{0em}{}

\setlist[enumerate]{itemsep=8pt}

\begin{document}

%--- Portada ---
\begin{center}
  {\Large\bfseries\color{tdblue} Universidad Torcuato Di Tella}\\[6pt]
  {\large\bfseries Tecnolog\'ia Digital VI}\\[4pt]
  {\large Banco de Preguntas Te\'oricas para el Parcial}\\[4pt]
  {\normalsize 1er Semestre 2026}\\[12pt]
  \rule{\linewidth}{1.5pt}
\end{center}

\begin{mdframed}[linecolor=tdorange, linewidth=1.5pt, roundcorner=4pt,
  backgroundcolor=lightyellow]
\textbf{C\'omo usar este documento:} Este banco cubre todos los temas te\'oricos
evaluables en el parcial. Las preguntas marcadas con
\textcolor{tdred}{\textbf{[CL\'ASICA]}} son las que se repiten hist\'oricamente.
Las respuestas aparecen en cajas verdes a continuaci\'on de cada pregunta.
\end{mdframed}

\vspace{0.5cm}
\tableofcontents
\newpage

% =============================================================
\section{Preprocesamiento y Validaci\'on}
% =============================================================

\subsection{Valores Faltantes}

\begin{enumerate}[label=\textbf{P\arabic*.}]

\item Explique los tres mecanismos de datos faltantes (MCAR, MAR, MNAR) y d\'e un ejemplo
  de cada uno.

\begin{answerbox}
\textbf{MCAR} (\textit{Missing Completely at Random}): la ausencia no depende de ning\'un
valor. Ejemplo: un sensor falla aleatoriamente.\\[4pt]
\textbf{MAR} (\textit{Missing at Random}): la ausencia depende de otras variables
observadas, pero no del valor faltante en s\'i. Ejemplo: hombres no responden preguntas
de salud con m\'as frecuencia que mujeres.\\[4pt]
\textbf{MNAR} (\textit{Missing Not at Random}): la ausencia depende del propio valor
faltante. Ejemplo: personas con ingresos muy altos no declaran su sueldo. Es el caso
m\'as problem\'atico.
\end{answerbox}

\item ¿Qu\'e estrategias existen para imputar valores faltantes? ¿Cu\'ando usar cada una?

\begin{answerbox}
\textbf{Media/mediana/moda}: simple, funciona bien con MCAR. La media es sensible a
outliers; la mediana es m\'as robusta.\\
\textbf{Nueva categor\'ia} (para categ\'oricas): modelar expl\'icitamente la ausencia.\\
\textbf{Imputaci\'on por modelo} (KNN, regresión): m\'as precisa pero costosa
computacionalmente.\\[4pt]
\textit{Importante}: toda imputaci\'on introduce alg\'un sesgo.
\end{answerbox}

\end{enumerate}

\subsection{Data Leakage}

\begin{enumerate}[label=\textbf{P\arabic*.}, resume]

\item \begin{classicbox}
\textcolor{tdred}{\textbf{[CL\'ASICA]}} ¿Qu\'e es el \textit{data leakage}? D\'e un
ejemplo concreto y explique c\'omo prevenirlo.
\end{classicbox}

\begin{answerbox}
El \textbf{data leakage} ocurre cuando informaci\'on del conjunto de test (o del futuro)
se filtra al entrenamiento. Consecuencia: m\'etricas artificialmente optimistas que no
se reproducen en producci\'on.\\[4pt]
\textbf{Ejemplos}:
\begin{itemize}[topsep=2pt, itemsep=2pt]
  \item Ajustar un StandardScaler sobre todo el dataset (train+test) antes de dividir.
  \item Imputar valores faltantes con la media de todo el dataset.
  \item Incluir variables que ``saben el futuro'' (ej: fecha de pago en un modelo de riesgo de impago).
\end{itemize}
\textbf{Prevenci\'on}: ajustar \emph{siempre} los transformadores solo con datos de
entrenamiento y luego aplicarlos al test. Usar pipelines.
\end{answerbox}

\item Un equipo aplica un MinMaxScaler ajustado sobre todo el dataset y luego divide en
  train/test. ¿Hay data leakage? Justifique.

\begin{answerbox}
\textbf{S\'i, hay data leakage.} El scaler aprendi\'o los m\'inimos y m\'aximos del
conjunto de test, de modo que el modelo ya ``vi\'o'' informaci\'on del test durante el
preprocesamiento. El flujo correcto es: dividir primero, luego ajustar el scaler solo
en train y transformar train y test con ese scaler.
\end{answerbox}

\end{enumerate}

\subsection{Escalado de Features}

\begin{enumerate}[label=\textbf{P\arabic*.}, resume]

\item ¿Cu\'ando es obligatorio escalar features y cu\'ando no importa? D\'e ejemplos de
  modelos en cada categor\'ia.

\begin{answerbox}
\textbf{Necesario}: modelos basados en distancias (KNN, SVM) o gradiente descendente
(regresi\'on log\'istica, redes neuronales). Sin escalado, features en escalas grandes
dominan artificialmente.\\[4pt]
\textbf{No necesario}: \'arboles de decisi\'on y sus ensambles (Random Forest, Gradient
Boosting), ya que sus splits comparan valores del mismo feature, independientemente de
la escala.
\end{answerbox}

\item Compare StandardScaler, MinMaxScaler y RobustScaler. ¿Cu\'ando preferir\'ia cada uno?

\begin{answerbox}
\textbf{StandardScaler}: centra en 0 con $\sigma=1$. Preferir cuando los datos son
aproximadamente normales y no hay outliers extremos.\\
\textbf{MinMaxScaler}: escala al rango $[0,1]$. \'Util cuando se necesita un rango
acotado, pero muy sensible a outliers.\\
\textbf{RobustScaler}: usa mediana e IQR. Preferir cuando hay outliers significativos,
ya que no se deja afectar por ellos.
\end{answerbox}

\end{enumerate}

\subsection{Desbalance de Clases}

\begin{enumerate}[label=\textbf{P\arabic*.}, resume]

\item \begin{classicbox}
\textcolor{tdred}{\textbf{[CL\'ASICA]}} Un dataset de detecci\'on de fraude tiene 98\%
de transacciones normales y 2\% fraudulentas. Un modelo que predice siempre ``normal''
obtiene 98\% de accuracy. ¿Qu\'e problema hay y qu\'e har\'ia?
\end{classicbox}

\begin{answerbox}
El accuracy es \textbf{enga\~noso} con clases desbalanceadas: un modelo trivial (siempre
predice la clase mayoritaria) ya alcanza 98\% sin aprender nada.\\[4pt]
\textbf{Soluciones}:
\begin{itemize}[topsep=2pt, itemsep=2pt]
  \item Usar m\'etricas adecuadas: F1-Score, Precision, Recall, AUC-PR.
  \item Oversampling de la clase minoritaria (SMOTE).
  \item Undersampling de la clase mayoritaria.
  \item Ajustar \texttt{class\_weight} en el modelo.
\end{itemize}
\end{answerbox}

\item ¿Qu\'e es SMOTE y c\'omo genera muestras sint\'eticas?

\begin{answerbox}
\textbf{SMOTE} (Synthetic Minority Oversampling Technique) genera muestras sint\'eticas
de la clase minoritaria interpolando entre muestras existentes y sus vecinos m\'as
cercanos en el espacio de features. No duplica muestras sino que crea nuevas
sint\'eticamente plausibles.
\end{answerbox}

\end{enumerate}

\subsection{Validaci\'on Cruzada}

\begin{enumerate}[label=\textbf{P\arabic*.}, resume]

\item \begin{classicbox}
\textcolor{tdred}{\textbf{[CL\'ASICA]}} Explique K-Fold Cross-Validation. ¿Por qu\'e
es preferible al simple holdout?
\end{classicbox}

\begin{answerbox}
En \textbf{K-Fold CV} se divide el dataset en $k$ particiones iguales. En cada iteraci\'on
se usa un fold como validaci\'on y los $k-1$ restantes como entrenamiento. Se promedian
las $k$ m\'etricas resultantes.\\[4pt]
\textbf{Ventajas sobre holdout simple}:
\begin{itemize}[topsep=2pt, itemsep=2pt]
  \item Cada muestra es usada exactamente una vez para validar.
  \item La estimaci\'on del error es m\'as robusta y con menor varianza.
  \item No depende de una \'unica divisi\'on aleatoria.
\end{itemize}
Valores t\'ipicos: $k=5$ o $k=10$.
\end{answerbox}

\item ¿Qu\'e es Leave-One-Out CV (LOOCV)? ¿Cu\'al es su ventaja y su desventaja principal?

\begin{answerbox}
\textbf{LOOCV} es K-Fold con $k=n$ (una muestra por fold). \\
\textbf{Ventaja}: usa casi todos los datos para entrenar en cada iteraci\'on.\\
\textbf{Desventaja}: es computacionalmente muy costoso ($n$ modelos a entrenar) y la
varianza de la estimaci\'on es alta.
\end{answerbox}

\end{enumerate}

% =============================================================
\section{M\'etricas de Evaluaci\'on}
% =============================================================

\subsection{Regresi\'on}

\begin{enumerate}[label=\textbf{P\arabic*.}, resume]

\item Compare MSE, RMSE y MAE. ¿Cu\'ando elegir\'ia cada uno?

\begin{answerbox}
\textbf{MSE}: penaliza cuadr\'aticamente los errores grandes. Usar cuando equivocarse
por mucho es muy costoso. Es diferenciable, est\'andar en entrenamiento.\\
\textbf{RMSE}: ra\'iz del MSE; mismas propiedades pero en la unidad original.
Comparar con MAE para detectar si hay puntos con error muy grande.\\
\textbf{MAE}: promedio de errores absolutos. Robusto a outliers (prioriza la mediana).
Usar cuando se quiere una m\'etrica directamente interpretable.
\end{answerbox}

\item ¿Puede el $R^2$ ser negativo? Explique qu\'e significar\'ia eso.

\begin{answerbox}
\textbf{S\'i}, el $R^2$ puede ser negativo. Significa que el modelo es \textbf{peor}
que simplemente predecir la media de los valores. $R^2 = 0$ equivale a predecir la
media; $R^2 = 1$ es predicci\'on perfecta.
\end{answerbox}

\item ¿Cu\'ando usar RMSLE en lugar de RMSE?

\begin{answerbox}
RMSLE es adecuado cuando importa el \textbf{error relativo (porcentual)} y no el absoluto.
Tambi\'en cuando los targets abarcan varios \'ordenes de magnitud. Penaliza m\'as las
subestimaciones que las sobreestimaciones. Muy usado en forecasting de ventas o precios.
\end{answerbox}

\end{enumerate}

\subsection{Clasificaci\'on}

\begin{enumerate}[label=\textbf{P\arabic*.}, resume]

\item \begin{classicbox}
\textcolor{tdred}{\textbf{[CL\'ASICA]}} Defina Precision y Recall. ¿Cu\'ando es m\'as
grave tener Recall bajo? ¿Cu\'ando es m\'as grave tener Precision baja?
\end{classicbox}

\begin{answerbox}
\[
  \text{Precision} = \frac{TP}{TP+FP}, \qquad \text{Recall} = \frac{TP}{TP+FN}
\]
\textbf{Recall bajo es m\'as grave} cuando los falsos negativos son catastr\'oficos:
detectar c\'ancer, fraude, fallas en reactores nucleares. ``No debo perder casos
positivos reales.''\\[4pt]
\textbf{Precision baja es m\'as grave} cuando los falsos positivos tienen alto costo:
filtrar spam (marcar emails leg\'itimos como spam), diagn\'osticos m\'edicos invasivos.
``No quiero molestar/perjudicar a quien no lo merece.''
\end{answerbox}

\item ¿Qu\'e es el F1-Score y cu\'ando es m\'as \'util que el Accuracy?

\begin{answerbox}
$F_1 = \dfrac{2 \cdot \text{Precision} \cdot \text{Recall}}{\text{Precision}+\text{Recall}}$

Es la media arm\'onica entre Precision y Recall. Es \'util cuando las clases est\'an
\textbf{desbalanceadas} y ambos tipos de error importan. El Accuracy puede ser enga\~noso
porque un clasificador trivial puede alcanzarlo sin aprender nada.
\end{answerbox}

\item Explique la curva ROC y el AUC-ROC. ¿Qu\'e significa un AUC de 0.5?

\begin{answerbox}
La \textbf{curva ROC} grafica TPR (Recall) vs.\ FPR para distintos umbrales de
clasificaci\'on. El \textbf{AUC-ROC} es el \'area bajo esa curva.\\
\begin{itemize}[topsep=2pt, itemsep=2pt]
  \item AUC = 1: clasificador perfecto.
  \item AUC = 0.5: equivalente a clasificar al azar (l\'inea diagonal).
  \item AUC < 0.5: peor que el azar (modelo invertido).
\end{itemize}
No usar ROC cuando las clases est\'an muy desbalanceadas; preferir la curva
Precision-Recall.
\end{answerbox}

\item ¿Cu\'ando preferir\'ia la curva Precision-Recall sobre la curva ROC?

\begin{answerbox}
Cuando el \textbf{dataset est\'a muy desbalanceado} y la clase positiva es la de inter\'es.
La curva ROC puede resultar artificialmente optimista porque el AUC se ve inflado por la
gran cantidad de verdaderos negativos (TN). La curva PR se focaliza solo en la clase
positiva y es m\'as informativa en estos casos.
\end{answerbox}

\end{enumerate}

% =============================================================
\section{Modelos Lineales}
% =============================================================

\subsection{Regresi\'on Lineal y Log\'istica}

\begin{enumerate}[label=\textbf{P\arabic*.}, resume]

\item ¿Por qu\'e la Regresi\'on Log\'istica es un modelo lineal si usa la funci\'on sigmoide?

\begin{answerbox}
La Regresi\'on Log\'istica es lineal porque la \textbf{frontera de decisi\'on} es un
hiperplano (combinaci\'on lineal de features = 0) y la relaci\'on entre features y
log-odds es lineal:
\[
  \log\!\left(\frac{p}{1-p}\right) = \beta_0 + \beta_1 x_1 + \cdots + \beta_p x_p
\]
La sigmoide solo transforma la salida lineal a probabilidades $(0,1)$, pero no cambia
la naturaleza lineal del modelo.
\end{answerbox}

\item ¿Por qu\'e la Regresi\'on Lineal tiene soluci\'on cerrada pero la Log\'istica no?

\begin{answerbox}
La Regresi\'on Lineal minimiza una funci\'on cuadr\'atica (MSE) que admite soluci\'on
exacta mediante las ecuaciones normales: $\hat{\boldsymbol{\beta}} = (\mathbf{X}^\top\mathbf{X})^{-1}\mathbf{X}^\top\mathbf{y}$.\\[4pt]
La Log\'istica minimiza la log-loss (cross-entrop\'ia), que no tiene soluci\'on cerrada
debido a la no linealidad de la sigmoide. Se requieren m\'etodos iterativos como
Gradient Descent.
\end{answerbox}

\end{enumerate}

\subsection{Regularizaci\'on}

\begin{enumerate}[label=\textbf{P\arabic*.}, resume]

\item \begin{classicbox}
\textcolor{tdred}{\textbf{[CL\'ASICA]}} Compare L1 (Lasso) y L2 (Ridge). ¿En qu\'e
se diferencian y cu\'ando usar cada una?
\end{classicbox}

\begin{answerbox}
\begin{center}
\begin{tabular}{lll}
\toprule
\textbf{Propiedad} & \textbf{L1 (Lasso)} & \textbf{L2 (Ridge)}\\
\midrule
Penalizaci\'on & $\lambda \sum |w_j|$ & $\lambda \sum w_j^2$\\
Coeficientes & Exactamente 0 (esparsidad) & Peque\~nos, nunca 0\\
Selecci\'on de features & S\'i & No\\
Multicolinealidad & Selecciona una & Distribuye peso\\
Soluci\'on & No anal\'itica & Anal\'itica\\
\bottomrule
\end{tabular}
\end{center}
\textbf{Usar L1} cuando se sospecha que muchas features son irrelevantes (selecci\'on
impl\'icita). \textbf{Usar L2} cuando todas las features son potencialmente \'utiles y
hay multicolinealidad.
\end{answerbox}

\item ¿Qu\'e efecto tiene aumentar $\lambda$ en la regularizaci\'on? ¿Y disminuirlo?

\begin{answerbox}
\textbf{Aumentar $\lambda$}: mayor penalizaci\'on sobre los coeficientes $\Rightarrow$
modelo m\'as simple $\Rightarrow$ mayor sesgo, menor varianza $\Rightarrow$ menos
overfitting.\\
\textbf{Disminuir $\lambda$}: menor penalizaci\'on $\Rightarrow$ modelo m\'as complejo
$\Rightarrow$ menor sesgo, mayor varianza $\Rightarrow$ m\'as riesgo de overfitting.\\
Con $\lambda=0$ no hay regularizaci\'on.
\end{answerbox}

\end{enumerate}

% =============================================================
\section{\'Arboles de Decisi\'on}
% =============================================================

\begin{enumerate}[label=\textbf{P\arabic*.}, resume]

\item Explique intuitivamente por qu\'e los \'arboles de decisi\'on son modelos de
  \textbf{alta varianza}.

\begin{answerbox}
Los \'arboles de decisi\'on son de alta varianza porque son muy sensibles a los datos
de entrenamiento: un cambio peque\~no en el dataset puede producir un \'arbol
completamente diferente. Esto ocurre porque el algoritmo CART es \textbf{goloso (greedy)}
y determinista: en cada nodo elige el mejor split para ese subconjunto, por lo que un
punto extra o diferente puede cambiar todos los splits subsiguientes en cascada.
\end{answerbox}

\item Defina Gini y Entrop\'ia. ¿Qu\'e valor toman cuando el nodo es puro? ¿Y cuando es
  m\'aximamente impuro?

\begin{answerbox}
\[
  \text{Gini}_m = 1 - \sum_{k=1}^{K} p_{mk}^2, \qquad
  H_m = -\sum_{k=1}^{K} p_{mk}\log_2(p_{mk})
\]
\textbf{Nodo puro} (todos de la misma clase): Gini = 0, Entrop\'ia = 0.\\
\textbf{M\'axima impureza} (clases equidistribuidas): Gini alcanza $1 - 1/K$ y
la Entrop\'ia alcanza $\log_2(K)$.
\end{answerbox}

\item \begin{classicbox}
\textcolor{tdred}{\textbf{[CL\'ASICA]}} Mencione \textbf{tres} hiperpar\'ametros que
regularizan un \'arbol de decisi\'on (reducen overfitting) y explique el efecto de cada
uno.
\end{classicbox}

\begin{answerbox}
\begin{enumerate}[label=\arabic*., topsep=2pt, itemsep=4pt]
  \item \texttt{max\_depth}: limita la profundidad m\'axima del \'arbol. Menos
    profundidad $\Rightarrow$ menos particiones $\Rightarrow$ menos overfitting.
  \item \texttt{min\_samples\_split}: m\'inimo de muestras requeridas para dividir
    un nodo. Un valor alto impide crear nodos con pocos datos.
  \item \texttt{min\_samples\_leaf}: m\'inimo de muestras en cada hoja. Evita hojas
    con una sola muestra que memorizan el ruido.
\end{enumerate}
(Tambi\'en v\'alido: \texttt{max\_features}, \texttt{max\_leaf\_nodes},
\texttt{ccp\_alpha} para poda.)
\end{answerbox}

\item Un \'arbol fue entrenado con temperaturas entre 10\textdegree C y 40\textdegree C.
  En producci\'on recibe una observaci\'on de 55\textdegree C. ¿Qu\'e predice? ¿Es un
  problema?

\begin{answerbox}
El \'arbol predecir\'a el valor de la \textbf{hoja m\'as extrema} para temperaturas
$> 40$\textdegree C (no puede extrapolar, solo puede asignar la \'ultima regi\'on).
\textbf{S\'i es un problema}: los \'arboles no extrapolan fuera del rango de
entrenamiento. Si los nuevos datos presentan valores fuera del rango visto en train, el
modelo tendr\'a un rendimiento degradado. Esto contrasta con los modelos lineales que
s\'i extrapolan (aunque tambi\'en con riesgo).
\end{answerbox}

\item ¿Qu\'e es la importancia de features en un \'arbol? ¿Cu\'al es su sesgo principal?

\begin{answerbox}
La \textbf{importancia de features} suma la ganancia en impureza ponderada que aporta
cada feature en todos los splits del \'arbol, normalizada para que sume 1.\\[4pt]
\textbf{Sesgo}: los \'arboles tienden a sobreestimar la importancia de variables con
alta cardinalidad o muchas categor\'ias, porque tienen m\'as oportunidades de hacer
buenos splits por azar. Correcci\'on: usar importancia por permutaci\'on.
\end{answerbox}

\end{enumerate}

% =============================================================
\section{Ensambles I: Bagging y Random Forest}
% =============================================================

\begin{enumerate}[label=\textbf{P\arabic*.}, resume]

\item Explique el concepto de \textbf{Bootstrap}. ¿Qu\'e porcentaje de muestras queda
  fuera (OOB) en promedio?

\begin{answerbox}
El \textbf{Bootstrap} consiste en extraer $n$ muestras \textbf{con reemplazo} de un
dataset de $n$ elementos. En promedio, $\approx 63.2\%$ de las muestras originales
aparecen al menos una vez y el $\approx 36.8\%$ (= $1/e$) queda fuera (\textit{Out-of-Bag},
OOB), ya que:
\[
  \lim_{n\to\infty}\left(1 - \frac{1}{n}\right)^n = \frac{1}{e} \approx 0.368
\]
\end{answerbox}

\item ¿C\'omo se usa el error OOB en Random Forest? ¿Qu\'e ventaja tiene frente a la
  validaci\'on cruzada?

\begin{answerbox}
Para cada muestra, el \textbf{error OOB} se estima usando solo los \'arboles que
\textbf{no} la incluyeron en su bootstrap. As\'i se obtiene una estimaci\'on del error
de generalizaci\'on sin necesidad de reservar un conjunto de validaci\'on separado.\\[4pt]
\textbf{Ventaja}: no desperdicia datos (todos se usan para entrenar) y es m\'as eficiente
computacionalmente que el K-Fold CV.
\end{answerbox}

\item \begin{classicbox}
\textcolor{tdred}{\textbf{[CL\'ASICA]}} ¿Cu\'al es la diferencia fundamental entre
Bagging y Random Forest?
\end{classicbox}

\begin{answerbox}
\textbf{Bagging} entrena m\'ultiples modelos en paralelo sobre muestras bootstrap
independientes. Reduce la varianza promediando predictores.\\[4pt]
\textbf{Random Forest} extiende Bagging con una capa adicional de aleatoriedad: en
cada split solo se eval\'uan $m$ features seleccionados al azar (t\'ipicamente
$m = \sqrt{p}$ para clasificaci\'on, $m = p/3$ para regresi\'on). Esto
\textbf{decorrelaciona m\'as} los \'arboles, reduciendo a\'un m\'as la varianza del
ensamble.
\end{answerbox}

\item ¿Agregar m\'as \'arboles a un Random Forest puede causar overfitting? ¿Y en
  Gradient Boosting?

\begin{answerbox}
\textbf{Random Forest}: \textbf{No}. Agregar m\'as \'arboles estabiliza o mejora el
rendimiento sin causar overfitting. El error OOB converge a un valor estable.\\[4pt]
\textbf{Gradient Boosting}: \textbf{S\'i}. Como los \'arboles se entrenan
secuencialmente para corregir los errores del anterior, agregar demasiados puede hacer
que el modelo se ajuste al ruido del entrenamiento (overfitting). Por eso se usa
\textit{early stopping}.
\end{answerbox}

\item ¿Por qu\'e la diversidad (decorrelaci\'on) entre modelos es importante en un
  ensamble?

\begin{answerbox}
Para $B$ modelos con varianza $\sigma^2$ y correlaci\'on $\rho$:
\[
  \text{Var}(\bar{f}) = \rho\,\sigma^2 + \frac{1-\rho}{B}\,\sigma^2
\]
Si $\rho \approx 0$ (modelos decorrelacionados), la varianza del ensamble se reduce a
$\sigma^2/B$, que tiende a 0 al aumentar $B$. Si $\rho$ es alto, el beneficio de
agregar modelos es m\'inimo. La clave es \textbf{reducir la correlaci\'on} entre
modelos.
\end{answerbox}

\end{enumerate}

% =============================================================
\section{Ensambles II: Boosting}
% =============================================================

\begin{enumerate}[label=\textbf{P\arabic*.}, resume]

\item \begin{classicbox}
\textcolor{tdred}{\textbf{[CL\'ASICA]}} Explique la diferencia conceptual entre Bagging
y Boosting en cuanto a c\'omo reducen el error.
\end{classicbox}

\begin{answerbox}
\textbf{Bagging}: entrena modelos en \textbf{paralelo} sobre muestras bootstrap
independientes. Reduce la \textbf{varianza} del modelo promediando predictores
diversos. No reduce mucho el sesgo.\\[4pt]
\textbf{Boosting}: entrena modelos \textbf{secuencialmente}, donde cada modelo nuevo
se focaliza en los errores del ensamble anterior. Reduce principalmente el \textbf{sesgo},
construyendo un modelo fuerte a partir de learners d\'ebiles.
\end{answerbox}

\item Explique el algoritmo de AdaBoost. ¿C\'omo se actualizan los pesos?

\begin{answerbox}
\begin{enumerate}[label=\arabic*., topsep=2pt, itemsep=2pt]
  \item Inicializar pesos iguales: $w_i = 1/n$.
  \item Entrenar clasificador $h_m$ con pesos $\{w_i\}$.
  \item Calcular error ponderado: $\varepsilon_m = \sum_{h_m(x_i)\neq y_i} w_i$.
  \item Peso del clasificador: $\alpha_m = \frac{1}{2}\ln\frac{1-\varepsilon_m}{\varepsilon_m}$.
  \item Aumentar pesos de muestras mal clasificadas; reducir las correctas.
  \item Renormalizar pesos.
\end{enumerate}
Predicci\'on final: $H(x) = \text{sign}\left(\sum_m \alpha_m h_m(x)\right)$.\\
Si $\varepsilon_m < 0.5$ entonces $\alpha_m > 0$ (clasificador mejor que el azar).
\end{answerbox}

\item ¿Qu\'e son los pseudo-residuos en Gradient Boosting? ¿Por qu\'e se usan?

\begin{answerbox}
Los \textbf{pseudo-residuos} son el gradiente negativo de la funci\'on de p\'erdida
respecto de la predicci\'on actual:
\[
  r_{im} = -\frac{\partial L(y_i, F(x_i))}{\partial F(x_i)}\bigg|_{F=F_{m-1}}
\]
Cada \'arbol nuevo se entrena para \textbf{predecir estos residuos}, lo que equivale
a dar un paso en la direcci\'on de mayor descenso de la p\'erdida en el espacio
funcional. Para MSE, el pseudo-residuo es simplemente $y_i - \hat{y}_i$.
\end{answerbox}

\item ¿Qu\'e rol cumplen el \textit{learning rate} ($\eta$) y el n\'umero de \'arboles
  ($M$) en Gradient Boosting? ¿C\'omo interact\'uan?

\begin{answerbox}
\textbf{Learning rate} $\eta$: controla cu\'anto contribuye cada \'arbol. Un $\eta$ bajo
hace pasos peque\~nos y m\'as conservadores.\\
\textbf{N\'umero de \'arboles} $M$: cu\'antos pasos se dan.\\[4pt]
\textbf{Trade-off}: $\eta$ peque\~no requiere $M$ grande para alcanzar buen desempe\~no,
pero produce mayor regularizaci\'on impl\'icita y mejor generalizaci\'on. Deben
ajustarse en conjunto.
\end{answerbox}

\item Compare XGBoost, LightGBM y CatBoost en al menos dos dimensiones.

\begin{answerbox}
\begin{center}
\begin{tabular}{llll}
\toprule
& \textbf{XGBoost} & \textbf{LightGBM} & \textbf{CatBoost}\\
\midrule
Crecimiento & Por nivel & Por hoja & Sim\'etrico\\
Velocidad & Moderada & Alta & Alta\\
Categ\'oricas & Requiere enc. & Nativo & Nativo\\
Regularizaci\'on & Expl\'icita & Impl.+expl. & Impl\'icita\\
\bottomrule
\end{tabular}
\end{center}
LightGBM es m\'as r\'apido pero m\'as propenso a overfitting por su crecimiento
\textit{leaf-wise}. CatBoost maneja categ\'oricas nativamente sin encoding previo.
\end{answerbox}

\end{enumerate}

% =============================================================
\section{Reducci\'on de la Dimensionalidad}
% =============================================================

\begin{enumerate}[label=\textbf{P\arabic*.}, resume]

\item ¿Qu\'e es la maldici\'on de la dimensionalidad y por qu\'e afecta a los modelos
  basados en distancias?

\begin{answerbox}
En espacios de alta dimensi\'on, el volumen crece exponencialmente. Los puntos se
vuelven \textbf{equidistantes} entre s\'i: la diferencia entre la distancia m\'axima y
m\'inima tiende a 0. Esto hace que los algoritmos basados en distancias (KNN, SVM radial,
clustering) pierdan efectividad porque la noci\'on de ``vecino cercano'' deja de ser
significativa.
\end{answerbox}

\item Explique PCA paso a paso. ¿Qu\'e maximiza y c\'omo se obtienen las componentes?

\begin{answerbox}
\textbf{PCA} busca una proyecci\'on lineal que \textbf{maximice la varianza retenida}.
\begin{enumerate}[label=\arabic*., topsep=2pt, itemsep=2pt]
  \item Estandarizar cada feature (media 0, $\sigma = 1$).
  \item Calcular la matriz de covarianza: $\mathbf{M}_x = \frac{1}{n-1}\mathbf{X}^\top\mathbf{X}$.
  \item Descomponer: $\mathbf{M}_x \mathbf{v}_j = \lambda_j \mathbf{v}_j$ (autovectores y autovalores).
  \item Ordenar por autovalor descendente.
  \item Proyectar: $\mathbf{Z} = \mathbf{X}\,\mathbf{P}_k$.
\end{enumerate}
La primera CP apunta a la direcci\'on de m\'axima varianza; las siguientes son
perpendiculares y tienen la segunda mayor varianza, etc.
\end{answerbox}

\item ¿Qu\'e significa la varianza explicada en PCA? ¿C\'omo se elige $k$?

\begin{answerbox}
La varianza explicada por la componente $j$ es:
\[
  \text{VE}_j = \frac{\lambda_j}{\sum_{l=1}^p \lambda_l}
\]
Para elegir $k$ (n\'umero de componentes a retener) se usa el \textbf{Scree Plot}: se
busca el $k$ donde la varianza acumulada supera un umbral deseado (t\'ipicamente
85--95\%), o donde la curva muestra un ``codo''.
\end{answerbox}

\item \begin{classicbox}
\textcolor{tdred}{\textbf{[CL\'ASICA]}} Compare PCA, t-SNE y UMAP en cuanto a:
linealidad, velocidad, y si pueden proyectar nuevas muestras.
\end{classicbox}

\begin{answerbox}
\begin{center}
\begin{tabular}{lccc}
\toprule
& \textbf{PCA} & \textbf{t-SNE} & \textbf{UMAP}\\
\midrule
Relaciones & Lineal & No lineal & No lineal\\
Velocidad & Muy r\'apida & Lenta & R\'apida\\
Estructura local & Media & Alta & Alta\\
Estructura global & Alta & Baja & Media\\
Proyecta nuevas muestras & S\'i & \textbf{No} & S\'i\\
Interpretabilidad & Alta & Baja & Baja\\
\bottomrule
\end{tabular}
\end{center}
t-SNE solo sirve para visualizaci\'on exploratoria; no puede proyectar nuevas muestras.
UMAP es preferido cuando se necesita velocidad y proyecci\'on de nuevas muestras.
\end{answerbox}

\end{enumerate}

% =============================================================
\section{Overfitting, Underfitting y Trade-offs}
% =============================================================

\begin{enumerate}[label=\textbf{P\arabic*.}, resume]

\item \begin{classicbox}
\textcolor{tdred}{\textbf{[CL\'ASICA --- SE TOMA SIEMPRE]}}
¿Qu\'e hiperpar\'ametros pueden causar \textbf{overfitting}? Liste al menos cinco
(de distintos modelos) y explique por qu\'e.
\end{classicbox}

\begin{answerbox}
\begin{enumerate}[label=\arabic*., topsep=2pt, itemsep=5pt]
  \item \textbf{\texttt{max\_depth} (\'arboles)}: valores altos permiten al \'arbol
    memorizar el ruido. Sin l\'imite, cada muestra puede tener su propia hoja.
  \item \textbf{\texttt{n\_estimators} (Gradient Boosting)}: demasiados \'arboles en
    boosting hacen que el modelo se ajuste al ruido de training (en RF no causa overfitting).
  \item \textbf{\texttt{learning\_rate} bajo con \texttt{n\_estimators} alto (GB)}: si la
    tasa es muy baja pero hay muchos \'arboles, el modelo puede sobreajustarse.
  \item \textbf{\texttt{min\_samples\_split} / \texttt{min\_samples\_leaf} bajos (\'arboles)}:
    valores peque\~nos permiten splits con muy pocos datos, creando ramas que memorizan
    outliers.
  \item \textbf{$\lambda$ de regularizaci\'on bajo (L1/L2)}: poca penalizaci\'on permite
    coeficientes muy grandes que se ajustan al ruido.
  \item \textbf{\texttt{max\_features} alto (RF)}: si se usan muchas features por split,
    los \'arboles quedan m\'as correlacionados y hay menos diversidad.
  \item \textbf{Grado del polinomio (regresi\'on polinomial)}: grados muy altos crean
    modelos que pasan exactamente por los datos de training.
  \item \textbf{$k$ peque\~no en KNN}: $k=1$ hace que el modelo simplemente memorice
    cada punto de entrenamiento.
\end{enumerate}
\end{answerbox}

\item Explique el Bias-Variance Trade-off con la f\'ormula y un ejemplo.

\begin{answerbox}
\[
  \mathbb{E}\!\left[(y - \hat{f}(x))^2\right]
  = \underbrace{\left(\mathbb{E}[\hat{f}(x)] - f(x)\right)^2}_{\text{Sesgo}^2}
  + \underbrace{\mathbb{E}\!\left[\left(\hat{f}(x) - \mathbb{E}[\hat{f}(x)]\right)^2\right]}_{\text{Varianza}}
  + \underbrace{\sigma^2_\varepsilon}_{\text{Ruido irred.}}
\]
\textbf{Ejemplo}: Un \'arbol muy profundo tiene bajo sesgo (se ajusta bien al train) pero
alta varianza (cambia mucho ante nuevos datos $\Rightarrow$ overfitting). Un modelo lineal
simple tiene alto sesgo (no captura patrones complejos $\Rightarrow$ underfitting) pero
baja varianza.
\end{answerbox}

\item ¿C\'omo detectar si un modelo tiene overfitting o underfitting a partir de las
  curvas de aprendizaje?

\begin{answerbox}
\textbf{Overfitting}: error de entrenamiento bajo pero error de validaci\'on alto (brecha
grande entre ambas curvas). El modelo memoriza el train pero no generaliza.\\[4pt]
\textbf{Underfitting}: error de entrenamiento \textbf{y} validaci\'on ambos altos (curvas
convergentes en valores altos). El modelo no captura la se\~nal de los datos.\\[4pt]
\textbf{Modelo adecuado}: error de validaci\'on tiene un m\'inimo; elegir la complejidad
en ese punto.
\end{answerbox}

\end{enumerate}

% =============================================================
\section{Selecci\'on de Features}
% =============================================================

\begin{enumerate}[label=\textbf{P\arabic*.}, resume]

\item Compare los tres tipos de selecci\'on de features: Filter, Wrapper y Embedded.

\begin{answerbox}
\begin{center}
\begin{tabular}{lll}
\toprule
\textbf{Tipo} & \textbf{M\'etodo} & \textbf{Caracter\'istica}\\
\midrule
Filter & Correlaci\'on, $\chi^2$, info.\ mutua & Independiente del modelo; r\'apido\\
Wrapper & Forward/Backward selection & Eval\'ua subconjuntos con modelo; lento\\
Embedded & L1 (Lasso), importancia en \'arboles & Integrado al entrenamiento\\
\bottomrule
\end{tabular}
\end{center}
\end{answerbox}

\item ¿Por qu\'e Label Encoding puede ser problem\'atico para modelos lineales?

\begin{answerbox}
El Label Encoding asigna un n\'umero entero a cada categor\'ia (ej: Argentina=0,
Brasil=1, Chile=2), lo que \textbf{impone un orden ordinal artificial} que el modelo
lineal interpretar\'a como relaci\'on num\'erica. El modelo entender\'a que
``Chile > Brasil > Argentina'', lo cual no tiene sentido si las categor\'ias no tienen
orden. Soluci\'on: usar One-Hot Encoding para variables nominales.\\
\textit{Excepci\'on}: en \'arboles de decisi\'on esto no es un problema, ya que los
splits comparan valores y no asumen relaciones lineales.
\end{answerbox}

\end{enumerate}

% =============================================================
\section{Resumen: Comparaci\'on de Modelos}
% =============================================================

\begin{enumerate}[label=\textbf{P\arabic*.}, resume]

\item \begin{classicbox}
\textcolor{tdred}{\textbf{[CL\'ASICA]}} ¿Cu\'ando elegiría un \'arbol simple sobre
Random Forest o Gradient Boosting?
\end{classicbox}

\begin{answerbox}
Se preferir\'ia un \'arbol simple cuando:
\begin{itemize}[topsep=2pt, itemsep=2pt]
  \item La \textbf{interpretabilidad} es cr\'itica (regulaci\'on, medicina, finanzas).
  \item Los recursos computacionales son muy limitados.
  \item El dataset es peque\~no y los ensambles podr\'ian no agregar valor.
  \item Se necesita un modelo de referencia (\textit{baseline}) r\'apido.
\end{itemize}
En todos los dem\'as casos, RF o GB suelen dar mejor performance.
\end{answerbox}

\item Complete la siguiente tabla comparativa de modelos:

\vspace{0.3cm}
\begin{center}
\begin{tabular}{lcccc}
\toprule
\textbf{Modelo} & \textbf{Sesgo} & \textbf{Varianza} & \textbf{Interpretable} & \textbf{No lineal}\\
\midrule
Reg.\ Lineal/Log\'istica & Alto & Bajo & Muy alta & No\\
\'Arbol (profundo) & Bajo & Alto & Media & S\'i\\
\'Arbol (podado) & Medio & Medio & Alta & S\'i\\
Bagging / RF & Bajo & Bajo & Baja & S\'i\\
Gradient Boosting & Muy bajo & Bajo & Baja & S\'i\\
\bottomrule
\end{tabular}
\end{center}

\end{enumerate}

% =============================================================
\section{Preguntas de Integraci\'on}
% =============================================================

\begin{enumerate}[label=\textbf{P\arabic*.}, resume]

\item Un equipo implementa el siguiente pipeline: (I) cargar dataset; (II) imputar con la
  media de todo el dataset; (III) escalar con MinMaxScaler sobre todo el dataset; (IV)
  dividir en train/test; (V) entrenar; (VI) evaluar. Identifique todos los errores.

\begin{answerbox}
\textbf{Dos errores de data leakage}:
\begin{enumerate}[label=\arabic*., topsep=2pt, itemsep=2pt]
  \item La imputaci\'on en el paso II usa la media del test: introduce info del futuro.
  \item El escalado en el paso III usa m\'inimos/m\'aximos del test.
\end{enumerate}
\textbf{Pipeline correcto}: (I) cargar; (II) dividir en train/test; (III) ajustar
imputador solo en train, aplicar a ambos; (IV) ajustar scaler solo en train, aplicar
a ambos; (V) entrenar; (VI) evaluar.
\end{answerbox}

\item Un modelo de detecci\'on de enfermedades tiene Recall = 0.99 pero Precision = 0.10.
  ¿C\'omo lo interpretar\'ia? ¿Es buen modelo?

\begin{answerbox}
\textbf{Recall = 0.99}: detecta el 99\% de los casos positivos reales (casi no pierde
enfermos). \textbf{Precision = 0.10}: de los que diagnostica como enfermos, solo el 10\%
realmente lo est\'a (90\% son falsos positivos).\\[4pt]
El modelo es \'util si el costo de no detectar una enfermedad es muy alto. Sin embargo,
la Precision baja genera muchos diagn\'osticos falsos que pueden causar tratamientos
innecesarios. Se deber\'ia ajustar el umbral o mejorar el modelo para balancear ambas
m\'etricas seg\'un el contexto cl\'inico.
\end{answerbox}

\item Un banco quiere predecir impago de cr\'editos. El dataset tiene 500k clientes,
  solo 12k no pagaron. ¿Qu\'e problemas anticipar\'ia y c\'omo los abordar\'ia?

\begin{answerbox}
\textbf{Problemas}:
\begin{enumerate}[label=\arabic*., topsep=2pt, itemsep=2pt]
  \item \textbf{Desbalance de clases} ($\approx 2.4\%$ positivos): el accuracy ser\'a
    enga\~noso. Usar F1, AUC-PR, ajustar \texttt{class\_weight}.
  \item \textbf{Valores faltantes} (8\% en ingreso, 3\% en historial): imputar
    correctamente, solo con datos de train.
  \item \textbf{Variables categ\'oricas} (empleo, zona, educaci\'on): aplicar encoding
    apropiado.
  \item \textbf{Data leakage potencial}: variables que dependen del outcome (ej: monto
    total de deuda acumulada post-impago).
\end{enumerate}
\textbf{M\'etrica}: AUC-PR o F1, no Accuracy.
\end{answerbox}

\end{enumerate}

\vspace{1cm}
\begin{center}
\rule{0.8\linewidth}{1pt}\\[6pt]
{\itshape Banco de preguntas completo --- TD VI 1S-2026}\\[4pt]
{\small\color{gray} Revise tambi\'en los parciales modelo 1, 2 y 3 para ejercicios con
gr\'aficos y matrices de confusi\'on.}
\end{center}

\end{document}
